\begin{lemma}
\label{lemma-nodal}
Let $k$ be a field. Let $X$ be a $1$-dimensional algebraic $k$-scheme.
Let $x \in X$ be a closed point. The following are equivalent
\begin{enumerate}
\item $x$ is a node,
\item $k \to \mathcal{O}_{X, x}$ is as in Lemma \ref{lemma-nodal-algebraic},
\item any $\overline{x} \in X_{\overline{k}}$ mapping to $x$ defines
a nodal singularity,
\item $\kappa(x)/k$ is separable, $\mathcal{O}_{X, x}$ is reduced, and
the first Fitting ideal of $\Omega_{X/k}$ generates $\mathfrak m_x$
in $\mathcal{O}_{X, x}$,
\item $\kappa(x)/k$ is separable, $\text{depth}(\mathcal{O}_{X, x}) = 1$, and
the first Fitting ideal of $\Omega_{X/k}$ generates $\mathfrak m_x$
in $\mathcal{O}_{X, x}$,
\item $\kappa(x)/k$ is separable and $\mathcal{O}_{X, x}$ is reduced, has
$\delta$-invariant $1$, and has $2$ geometric branches.
\end{enumerate}
\end{lemma}

\begin{proof}
First assume that $k$ is algebraically closed.
In this case the equivalence of (1) and (3) is trivial.
The equivalence of (1) and (3) with (2) holds because the only
nondegenerate quadric in two variables is $xy$ up to change in
coordinates. The equivalence of (1) and (6) is 
Lemma \ref{lemma-multicross-algebra}.
After replacing $X$ by an affine neighbourhood
of $x$, we may assume there is a closed immersion $X \to \mathbf{A}^n_k$
mapping $x$ to $0$. Let $f_1, \ldots, f_m \in k[x_1, \ldots, x_n]$
be generators for the ideal $I$ of $X$ in $\mathbf{A}^n_k$.
Then $\Omega_{X/k}$ corresponds to the $R = k[x_1, \ldots, x_n]/I$-module
$\Omega_{R/k}$ which has a presentation
$$
R^{\oplus m} \xrightarrow{(\partial f_j/\partial x_i)}
R^{\oplus n} \to \Omega_{R/k} \to 0
$$
(See Algebra, Sections \ref{algebra-section-differentials} and
\ref{algebra-section-netherlander}.)
The first Fitting ideal of $\Omega_{R/k}$ is thus the ideal
generated by the $(n - 1) \times (n - 1)$-minors of the
matrix $(\partial f_j/\partial x_i)$. Hence (2), (4), (5)
are equivalent by Lemma \ref{lemma-fitting-ideal} applied
to the completion of $k[x_1, \ldots, x_n] \to R$
at the maximal ideal $(x_1, \ldots, x_n)$.

\medskip\noindent
Now assume $k$ is an arbitrary field.
In cases (2), (4), (5), (6) the residue field $\kappa(x)$ is
separable over $k$. Let us show this holds as well
in cases (1) and (3). Namely, let $Z \subset X$ be the closed subscheme
of $X$ defined by the first Fitting ideal of $\Omega_{X/k}$.
The formation of $Z$ commutes with field extension
(Divisors, Lemma \ref{divisors-lemma-base-change-and-fitting-ideal-omega}).
If (1) or (3) is true, then there exists a point
$\overline{x}$ of $X_{\overline{k}}$ such that $\overline{x}$
is an isolated point of multiplicity $1$ of $Z_{\overline{k}}$ (as we have the
equivalence of the conditions of the lemma over $\overline{k}$).
In particular $Z_{\overline{x}}$ is geometrically reduced at $\overline{x}$
(because $\overline{k}$ is algebraically closed). Hence
$Z$ is geometrically reduced at $x$
(Varieties, Lemma \ref{varieties-lemma-geometrically-reduced-upstairs}).
In particular, $Z$ is reduced at $x$, hence $Z = \Spec(\kappa(x))$
in a neighbourhood of $x$ and $\kappa(x)$ is geometrically reduced
over $k$. This means that $\kappa(x)/k$ is separable
(Algebra, Lemma \ref{algebra-lemma-characterize-separable-field-extensions}).

\medskip\noindent
The argument of the previous paragraph shows that if (1) or (3) holds, then
the first Fitting ideal of $\Omega_{X/k}$ generates $\mathfrak m_x$.
Since $\mathcal{O}_{X, x} \to \mathcal{O}_{X_{\overline{k}}, \overline{x}}$
is flat and since $\mathcal{O}_{X_{\overline{k}}, \overline{x}}$
is reduced and has depth $1$, we see that (4) and (5) hold
(use Algebra, Lemmas \ref{algebra-lemma-descent-reduced} and
\ref{algebra-lemma-apply-grothendieck}).
Conversely, (4) implies (5) by
Algebra, Lemma \ref{algebra-lemma-criterion-reduced}.
If (5) holds, then $Z$ is geometrically reduced at $x$
(because $\kappa(x)/k$ separable and $Z$ is $x$ in a neighbourhood).
Hence $Z_{\overline{k}}$ is reduced at any point $\overline{x}$
of $X_{\overline{k}}$ lying over $x$. In other words, the first
fitting ideal of $\Omega_{X_{\overline{k}}/\overline{k}}$ generates
$\mathfrak m_{\overline{x}}$ in $\mathcal{O}_{X_{\overline{k}, \overline{x}}}$.
Moreover, since
$\mathcal{O}_{X, x} \to \mathcal{O}_{X_{\overline{k}}, \overline{x}}$ is flat
we see that $\text{depth}(\mathcal{O}_{X_{\overline{k}}, \overline{x}}) = 1$
(see reference above).
Hence (5) holds for $\overline{x} \in X_{\overline{k}}$ and we
conclude that (3) holds (because of the equivalence over algebraically
closed fields). In this way we see that (1), (3), (4), (5)
are equivalent.

\medskip\noindent
The equivalence of (2) and (6) follows from
Lemma \ref{lemma-2-branches-delta-1}.

\medskip\noindent
Finally, we prove the equivalence of (2) $=$ (6) with
(1) $=$ (3) $=$ (4) $=$ (5). First we note that the geometric number
of branches of $X$ at $x$ and the geometric number of branches
of $X_{\overline{k}}$ at $\overline{x}$ are equal by
Varieties, Lemma
\ref{varieties-lemma-geometric-branches-and-change-of-fields}.
We conclude from the information available to us at this point
that in all cases this number is equal to $2$.
On the other hand, in case (1) it is clear that $X$ is geometrically
reduced at $x$, and hence
$$
\delta\text{-invariant of }X\text{ at }x \leq
\delta\text{-invariant of }X_{\overline{k}}\text{ at }\overline{x}
$$
by Varieties, Lemma \ref{varieties-lemma-delta-invariant-and-change-of-fields}.
Since in case (1) the right hand side is $1$, this
forces the $\delta$-invariant of $X$ at $x$ to be $1$
(because if it were zero, then $\mathcal{O}_{X, x}$ would
be a discrete valuation ring by
Varieties, Lemma \ref{varieties-lemma-delta-invariant-is-zero}
which is unibranch, a contradiction). Thus (5) holds.
Conversely, if (2) $=$ (5) is true, then assumptions (a), (b), (c) of
Varieties, Lemma \ref{varieties-lemma-geometrically-normal-in-codim-1}
hold for $x \in X$ by
Lemma \ref{lemma-2-branches-delta-1}. Thus
Varieties, Lemma
\ref{varieties-lemma-delta-invariant-and-change-of-fields-better}
applies and shows that we have equality in the above displayed inequality.
We conclude that (5) holds for $\overline{x} \in X_{\overline{k}}$
and we are back in case (1) by the equivalence of the conditions
over an algebraically closed field.
\end{proof}